\subsection{Сутнісна характеристика та структура процесу навчання розробки імерсивних освітніх ресурсів}\label{subsec:1.2}

Трансформації освітнього середовища, зумовлені цифровізацією, поширенням технологій занурення та зміною ролі вчителя в умовах Нової української школи, актуалізують переосмислення процесу професійної підготовки майбутніх учителів інформатики. Аналітичне вивчення наукової літератури (підрозділ~\ref{subsec:1.1}) свідчить про те, що процес навчання розробки імерсивних освітніх ресурсів доцільно розглядати не лише як сукупність навчальних дій і технологічних операцій, а як цілісну систему, що інтегрує мотиваційні, когнітивні, діяльнісні та рефлексивні компоненти професійної підготовки майбутнього вчителя. Сутність процесу визначається взаємодією педагогічних, психологічних і технологічних чинників, а також орієнтацією майбутнього вчителя на досягнення освітньо значущих результатів, пов'язаних із проєктуванням безпечних, доступних і педагогічно доцільних імерсивних ресурсів.

З огляду на викладене виникає потреба в окресленні структурної організації означеного процесу, визначенні його ключових складників, логіки їхньої взаємодії та функціонального призначення в системі професійної підготовки майбутніх учителів інформатики. Передусім доцільно окреслити термінологічний апарат дослідження, оскільки процес навчання розробки імерсивних освітніх ресурсів пов'язаний із використанням значної кількості міждисциплінарних понять, що функціонують у педагогічному, психологічному та технологічному дискурсах і не завжди мають усталене трактування в межах професійної педагогічної освіти.

У межах дослідження суттєвого значення набувають поняття, що безпосередньо відображають специфіку імерсивних технологій як об'єкта педагогічного проєктування, роль цифрового дизайну в освітньому процесі та професійну підготовку майбутніх учителів інформатики. До таких базових понять належать: ``імерсивні технології'' (далі -- ІТ) як сукупність цифрових засобів і середовищ занурення; ``імерсивні освітні ресурси'' (далі -- ІОР) як педагогічно спроєктовані цифрові продукти; процес навчання розробки імерсивних освітніх ресурсів (далі -- ПНР~ІОР); цифровий дизайн (далі -- ЦД) в освіті та його диференційовані варіанти (digital instructional design, digital learning design, digital media design); професійна готовність майбутнього вчителя інформатики до розробки імерсивних освітніх ресурсів; педагогічна ефективність імерсивних освітніх ресурсів. Окреслені поняття потребують детальнішого розгляду з метою конкретизації їх змісту та ролі у структурі досліджуваного процесу.

Водночас окреслення термінологічного апарату потребує опори на міжнародні теоретичні моделі, що розкривають механізми навчання в імерсивних середовищах та визначають аналітичну рамку подальшого розгляду. Фундаментальною для дослідження є модель навчальних афордансів тривимірних віртуальних середовищ Б.~Далґарно та М.~Дж.~В.~Лі~\cite{Dalgarno2010What}, яка ідентифікує два ключових параметри -- репрезентаційну точність (\emph{representational fidelity}) та взаємодію учня (\emph{learner interaction}) -- що спільно породжують п'ять навчальних афордансів: просторові знання, експерієнціальне навчання, залученість і мотивацію, контекстуальне навчання та колаборативне навчання. Ця модель обґрунтовує, що дидактичний потенціал імерсивних технологій визначається не рівнем технологічного занурення як таким, а взаємодією між якістю цифрового представлення та інтерактивними можливостями середовища.

К.~Фаулер~\cite{Fowler2015Pedagogy} поглиблює цю модель, вводячи поняття \emph{педагогічного занурення}: три стадії навчання (концептуалізація, конструювання, діалог) відповідають трьом типам використання віртуальних середовищ. Ключовий висновок полягає у тому, що занурення є не лише технологічною властивістю середовища, а педагогічним проєктувальним рішенням, яке вчитель має усвідомлено обирати відповідно до навчальних цілей. М.~Мюлдерс, Й.~Бухнер та М.~Керрес~\cite{Mulders2020Framework} конкретизують ці положення для імерсивної віртуальної реальності, обґрунтовуючи на основі когнітивної теорії мультимедійного навчання (\emph{CTML}), що унікальні характеристики іVR (висока присутність, тілесна взаємодія, просторовий контекст) вимагають модифікованих принципів педагогічного дизайну, які не зводяться до прямого перенесення мультимедійних рекомендацій.

У сукупності з когнітивно-афективною моделлю імерсивного навчання CAMIL~\cite{Makransky2021CAMIL}, проаналізованою у підрозділі~\ref{subsec:1.1}, зазначені теоретичні положення формують аналітичну рамку для подальшого визначення сутнісних характеристик імерсивних освітніх ресурсів: кожна характеристика ІОР має бути обґрунтована не лише технологічно, а й з позицій навчальних афордансів, педагогічного занурення та когнітивних процесів учасників.

Передусім зазначимо, що базове поняття ``імерсивність'' використовується у різних аспектах та переважно визначається як занурення, тобто стан глибокої когнітивної, емоційної та сенсорної залученості суб'єкта в певну діяльність або штучно створене середовище, що забезпечує ефект присутності та інтенсифікує взаємодію з навчальним або віртуальним простором~\cite{Semerikov2023Presence}.

\textbf{Імерсивні технології} у сучасному науково-педагогічному дискурсі розглядаються як сукупність цифрових засобів і середовищ, що забезпечують повне або часткове занурення суб'єкта навчання у віртуально сконструйований простір. У науковій літературі зазначені технології тлумачаться як інструменти створення ефекту присутності та взаємодії людини з цифровим середовищем, що передбачає її залучення до штучно змодельованого або змішаного простору \cite{Semerikov2023Presence} Такий підхід поділяють і зарубіжні науковці, зокрема Й.~Радіанті (J. Radianti) та співавтори, Дж.~Паронг (J. Parong) і Р.~Е.~Маєр (R. E. Mayer), Г.~Макранскі (G. Makransky), у працях яких імерсивність (\emph{immersion}) трактується передусім як властивість технологічного середовища навчання, зумовлена рівнем технічного занурення та формуванням відчуття присутності \emph{(presence)}, а не як самостійно концептуалізований освітній ресурс \cite{Parong2021Cognitive}. Водночас у зарубіжному науковому дискурсі представлена й інша дослідницька позиція, відповідно до якої імерсивність не зводиться виключно до технічних характеристик середовища або використання спеціалізованих пристроїв занурення. Так, у працях Дж.~Г.~Робертсона (G. G. Robertson), М.~Червінські (M. Czerwinski) та М.~ван Данціха (M. van Dantzich) імерсивність розглядається як психологічний ефект присутності, що може формуватися і в настільних (\emph{desktop}) віртуальних середовищах за умови відповідного дизайну взаємодії, інтерфейсу та організації уваги користувача. У такому підході імерсивність постає не як похідна від рівня сенсорного занурення, а як результат цілісної взаємодії користувача з цифровим середовищем, що розширює межі розуміння імерсивних технологій поза рамками виключно повністю імерсивної віртуальної реальності\cite{Krokos2019Virtual}.

Також звертаємо увагу на те, що в зарубіжному науковому дискурсі імерсивність та присутність не ототожнюються автоматично з використанням повністю імерсивних пристроїв, а розглядаються крізь призму когнітивного та психологічного досвіду користувача. Систематичний огляд досліджень, здійснений Д.~Гамільтоном (D.~Hamilton), Дж.~Маккечні (J.~McKechnie), Е.~Едгелл (E.~Edgell) та К.~Мак (C.~Mack) \cite{Kennedy2023Immersive}, засвідчує, що імерсивна віртуальна реальність як педагогічний інструмент здатна посилювати залученість здобувачів освіти та сприяти глибшому засвоєнню навчального матеріалу завдяки створенню ефекту присутності та активній взаємодії з навчальним контентом. Автори наголошують на необхідності педагогічно обґрунтованого проєктування імерсивних навчальних середовищ, оскільки саме дидактичний дизайн, а не рівень технологічного занурення, визначає ефективність імерсивної віртуальної реальності як засобу навчання.

У розширеному трактуванні імерсивні технології охоплюють різні форми віртуальної, доповненої та змішаної реальності, які забезпечують як повне занурення у віртуальний світ, так і інтеграцію цифрових об'єктів у реальне середовище~\cite{DiNatale2020Immersive}. Зважаючи на існування реальної (об'єктивної) реальності \emph{(RR~-- real reality)}, в якій знаходиться користувач і яку він сприймає органами своїх чуттів \cite{Omelchuk2024Secondary}, на думку вчених, головна ідея імерсивних технологій полягає у підвищенні відчуття присутності, взаємодії та участі користувача у віртуальному або поєднаному (реальному та віртуальному) середовищі \cite{Omelchuk2024Secondary, Semerikov2021Ref}.

Відповідно до класичного визначення Р.~Азуми~\cite{TechGlossaryAR2020}, \emph{доповнена реальність} (AR) характеризується поєднанням реальних і віртуальних об'єктів, інтерактивністю в режимі реального часу та тривимірною реєстрацією. У таксономії континууму ``реальність--віртуальність'' П.~Мілграма та Ф.~Кішіно~\cite{KVR2017Definition} \emph{віртуальна реальність} (VR) розміщується на протилежному від реальності полюсі як повністю синтезоване середовище. \emph{Змішана реальність} (MR) займає проміжне положення, забезпечуючи інтерактивну інтеграцію цифрових об'єктів з фізичним простором. Принциповою для цього дослідження є також \emph{веб-доповнена реальність} (WebAR)~-- реалізація AR-досвіду безпосередньо у веб-браузері без встановлення додатків, що забезпечує найвищий рівень доступності та релевантна для підготовки майбутніх учителів інформатики завдяки використанню відкритих веб-технологій (Three.js, MindAR, WebXR API).

Систематизоване порівняння основних різновидів імерсивних технологій за ключовими параметрами подано у табл.~\ref{tab:ar_vr_mr_comparison}.

\begin{table}[!h]
\caption{Порівняльна характеристика основних різновидів імерсивних технологій}\label{tab:ar_vr_mr_comparison}\small
\begin{tabularx}{\linewidth}{|p{0.18\linewidth}|p{0.19\linewidth}|p{0.17\linewidth}|p{0.19\linewidth}|X|}
\hline
\hfil \textbf{Параметр} & \hfil \textbf{AR} & \hfil \textbf{VR} & \hfil \textbf{MR} & \hfil \textbf{WebAR} \\
\hline
Середовище & Реальне + цифрове накладання & Повністю синтетичне & Реальне + інтерактивне цифрове & Реальне + браузерне цифрове \\
\hline
Рівень занурення & Часткове & Повне & Гібридне & Часткове (мобільне) \\
\hline
Обладнання & Смартфон / планшет & VR-шолом & MR-окуляри & Веб-браузер \\
\hline
Модель взаємодії & Маркер / GPS / зображення & Контролер / жести & Просторове картування & Дотик / камера \\
\hline
Доступність & Висока & Низька & Низька & Найвища (без інсталяції) \\
\hline
Фреймворки & ARKit, ARCore & Unity, Unreal & HoloLens SDK & MindAR, Three.js, WebXR API \\
\hline
\end{tabularx}
\end{table}


Як засвідчує табл.~\ref{tab:ar_vr_mr_comparison}, WebAR поєднує переваги доповненої реальності з доступністю веб-технологій, що робить її оптимальним інструментом для підготовки майбутніх учителів інформатики: відсутність потреби у спеціалізованому обладнанні, використання відкритих стандартів і можливість розгортання на будь-якому пристрої з веб-браузером. Окрему групу імерсивних технологій становлять \emph{360° відео й фотографії}, що формують панорамний ефект присутності, а також \emph{інтерактивні симуляції,} спрямовані на моделювання реальних або умовних ситуацій з можливістю прийняття рішень і спостереження за їх наслідками.

\begin{figure}[!h]
\centering
%\resizebox{\textwidth}{!}{%
\begin{tikzpicture}

% Title
\node[dia/header=10cm] (title) {ІМЕРСИВНІ ТЕХНОЛОГІЇ};

% Row 1
\node[dia/rbox=4.5cm, below=1.5cm of title.south, xshift=-4cm] (l1)
  {Віртуальна реальність};
\node[dia/rbox=6cm, right=1cm of l1] (r1)
  {Спеціалізована технічні засоби};

% Row 2
\node[dia/rbox=4.5cm, below=0.8cm of l1] (l2)
  {Доповнена реальність};
\node[dia/rbox=6cm, right=1cm of l2] (r2)
  {Цифрова трансформація + реальний простір};

% Row 3
\node[dia/rbox=4.5cm, below=0.8cm of l2] (l3)
  {Змішана реальність};
\node[dia/rbox=6cm, right=1cm of l3] (r3)
  {Інтеграція цифрових об'єктів у реальний простір};

% Row 4
\node[dia/rbox=4.5cm, below=0.8cm of l3] (l4)
  {360° відео й фотографії};
\node[dia/rbox=6cm, right=1cm of l4] (r4)
  {Панорамний ефект присутності};

% Row 5
\node[dia/rbox=4.5cm, below=0.8cm of l4] (l5)
  {Інтерактивні симуляції};
\node[dia/rbox=6cm, right=1cm of l5] (r5)
  {Моделювання ситуацій};

% Horizontal arrows from left to right boxes
\draw[dia/arrow] (l1.east) -- (r1.west);
\draw[dia/arrow] (l2.east) -- (r2.west);
\draw[dia/arrow] (l3.east) -- (r3.west);
\draw[dia/arrow] (l4.east) -- (r4.west);
\draw[dia/arrow] (l5.east) -- (r5.west);

% Vertical connector: from title left side, straight down, with horizontal ticks
\coordinate (bar x) at ([xshift=-0.5cm]l1.west);
% Vertical line from title down to last box level
\draw[dia/line] ([xshift=-2.5cm]title.south) -- (bar x |- title.south)
  -- (bar x |- l5.west);
% Horizontal ticks from vertical bar to each left box
\foreach \n in {l1,l2,l3,l4,l5} {
  \draw[dia/line] (bar x |- \n.west) -- (\n.west);
}

\end{tikzpicture}%
%}

\caption{Основні різновиди імерсивних технологій}
\label{fig:1.3}
\end{figure}

\emph{\textbf{Імерсивні освітні ресурси}} постають як результат педагогічно вмотивованого застосування імерсивних технологій у освітньому процесі та розглядаються як спеціально спроєктовані цифрові продукти, спрямовані на досягнення визначених освітніх цілей. На відміну від традиційних навчальних матеріалів, такі ресурси ґрунтуються на використанні віртуальної, доповненої, змішаної реальності та інших форм імерсивних середовищ, що забезпечують занурення здобувачів освіти в навчальну діяльність, активну взаємодію з контентом і формування практично орієнтованих умінь у змодельованих або наближених до реальних умовах.

Спираючись на аналітичну рамку, обґрунтовану вище, виокремимо п'ять сутнісних характеристик ІОР з їх теоретичним обґрунтуванням:

\begin{enumerate}[label=\arabic*.]
	\item 
\emph{Інтерактивність} -- активна участь користувача у взаємодії з навчальним середовищем і вплив на перебіг навчальної діяльності. У моделі Б.~Далґарно та М.~Дж.~В.~Лі~\cite{Dalgarno2010What} взаємодія учня (\emph{learner interaction}) виступає одним із двох фундаментальних параметрів, що визначають навчальні афорданси 3D-середовищ: швидкість відгуку, діапазон можливих дій та відповідність між намірами користувача і реакціями системи безпосередньо впливають на якість навчального досвіду.

	\item 
\emph{Занурення / присутність} -- ефект перебування у віртуальному або змішаному освітньому просторі, що сприяє глибшому засвоєнню навчального матеріалу. Принципово важливим є розрізнення занурення як технологічної властивості середовища та присутності як психологічного стану~\cite{Semerikov2023Presence}. У моделі CAMIL~\cite{Makransky2021CAMIL} присутність є першим із шести факторів, що опосередковують вплив імерсивних технологій на навчальні результати. К.~Фаулер~\cite{Fowler2015Pedagogy} уточнює, що педагогічне занурення залежить від свідомого дизайнерського рішення вчителя, а не лише від технічних параметрів пристрою.

	\item 
\emph{Контекстуальність} -- навчання у межах реальних або спеціально змодельованих ситуацій, що підвищує практичну значущість набутих знань. Ця характеристика відповідає афордансу ``контекстуального навчання'' у моделі Далґарно та Лі, а також узгоджується з положеннями ситуативного пізнання, за якими знання формується не ізольовано, а у процесі автентичної діяльності в контексті його подальшого застосування.

	\item 
\emph{Візуалізація} -- використання цифрових візуальних засобів для представлення складних понять і процесів. У концепції Далґарно та Лі~\cite{Dalgarno2010What} репрезентаційна точність (\emph{representational fidelity}) визначається як другий ключовий параметр 3D-середовищ, а М.~Мюлдерс, Й.~Бухнер та М.~Керрес~\cite{Mulders2020Framework} обґрунтовують, що для іVR цей параметр вимагає дотримання модифікованих принципів когнітивної теорії мультимедійного навчання.

	\item 
\emph{Адаптивність} -- можливість варіативного налаштування змісту й складності ресурсу відповідно до індивідуальних освітніх потреб. Обґрунтування цієї характеристики спирається на теорію когнітивного навантаження: М.~Пупар зі співавторами~\cite{Poupard2024Systematic} встановили, що доповнена реальність оптимізує когнітивне навантаження для початківців, тоді як VR може його збільшувати, що зумовлює потребу в адаптивному управлінні складністю ІОР.
\end{enumerate}

Взаємозв'язок п'яти сутнісних характеристик ІОР та їх теоретичного підґрунтя представлено на рис.~\ref{fig:ier_chars}.

\begin{figure}[!h]
\centering
%\resizebox{\textwidth}{!}{%
\begin{tikzpicture}

% Central node
\node[dia/rbox=4cm, minimum height=1.5cm] (center) at (0,0)
  {\textbf{Педагогічна \\ефективність ІОР}};

% Pentagon vertices (5 characteristics with theoretical anchors)
% Top
\node[dia/rbox=4cm, minimum height=1.8cm] (int) at (90:5cm)
  {\textbf{Інтерактивність}\\[2pt]\scriptsize Dalgarno \& Lee:\\взаємодія учня};

% Upper-left
\node[dia/rbox=3.5cm, minimum height=1.8cm] (imm) at (162:5cm)
  {\textbf{Занурення / присутність}\\[2pt]\scriptsize Slater; CAMIL:\\психологічна присутність};

% Lower-left
\node[dia/rbox=4.5cm, minimum height=1.8cm] (ctx) at (234:5cm)
  {\textbf{Контекстуальність}\\[2pt]\scriptsize Ситуативне пізнання;\\експерієнціальне навчання};

% Lower-right
\node[dia/rbox=4.5cm, minimum height=1.8cm] (vis) at (306:5cm)
  {\textbf{Візуалізація}\\[2pt]\scriptsize CTML; Mulders et al.:\\репрезентаційна точність};

% Upper-right
\node[dia/rbox=3.5cm, minimum height=1.8cm] (adp) at (18:5cm)
  {\textbf{Адаптивність}\\[2pt]\scriptsize CLT; Poupard et al.:\\когнітивне навантаження};

% Arrows from vertices to center
\draw[dia/arrow] (int) -- (center);
\draw[dia/arrow] (imm) -- (center);
\draw[dia/arrow] (ctx) -- (center);
\draw[dia/arrow] (vis) -- (center);
\draw[dia/arrow] (adp) -- (center);

% Pentagon edges (dashed)
\draw[dia/line, dashed] (int) -- (imm);
\draw[dia/line, dashed] (imm) -- (ctx);
\draw[dia/line, dashed] (ctx) -- (vis);
\draw[dia/line, dashed] (vis) -- (adp);
\draw[dia/line, dashed] (adp) -- (int);

\end{tikzpicture}%
%}

\caption{Сутнісні характеристики імерсивних освітніх ресурсів та їх теоретичне підґрунтя}
\label{fig:ier_chars}
\end{figure}

У межах дослідження було здійснено класифікацію імерсивних освітніх ресурсів з урахуванням їхнього функціонального призначення, дидактичної ролі та ступеня безпосередньої залученості до освітнього процесу. Тобто ІОР класифікуються не за формальною ознакою їхнього типу, а за дидактичною функцією в процесі навчання розробки імерсивних освітніх ресурсів майбутніх учителів інформатики. ІОР вже не лише не як сукупність цифрових продуктів, а більше - інструменти організації навчальної діяльності майбутніх учителів інформатики. Відповідно до запропонованого підходу імерсивні освітні ресурси доцільно поділяти на основні та допоміжні, що відображає їхню різну роль у навчанні розробки імерсивних освітніх ресурсів.

До основних ІОР віднесено імерсивні навчально-методичні комплекси, системи підтримки навчання, навчальні лабораторії та предметні середовища, які безпосередньо забезпечують формування професійних умінь студентів і створюють умови для інтерактивної взаємодії з віртуальними або змішаними середовищами. У межах цієї групи виокремлюються імерсивні програмно-методичні матеріали, засоби оцінювання навчальних досягнень, практикуми, навчально-методичні матеріали та середовища моделювання, що реалізують різні дидактичні функції - від подання навчального контенту до організації симуляційної та проєктної діяльності.

Допоміжні ІОР представлені додатковими науково-навчальними матеріалами та комунікаційними засобами, які розширюють можливості основних ІОР, забезпечують інформаційну підтримку, зворотний зв'язок і координацію навчальної взаємодії. Структуризація ІОР, представлена на рис.~\ref{fig:1.4}, спирається на результати попередньої апробації авторського підходу, здійсненої в межах кваліфікаційного дослідження~\cite{Shepiliev2019Master}, та адаптованої до завдань підготовки майбутніх учителів інформатики.

\begin{figure}[!h]
\centering
%\resizebox{\textwidth}{!}{%
\begin{tikzpicture}

% Title
\node[dia/header=14cm] (title) {ІМЕРСИВНІ ОСВІТНІ РЕСУРСИ};

% --- Upper-left --- use xshift for center-to-center offset
\node[dia/rbox=5.5cm, below=1.5cm of title.south, xshift=-4.5cm] (h1)
  {ІОР формування\\ проєктно-конструкторських умінь};
\node[dia/lbox=5.5cm, below=0.3cm of h1] (d1) {%
  -- середовища моделювання\\
  -- навчальні лабораторії\\
  -- предметні середовища%
};

% --- Upper-right ---
\node[dia/rbox=5.5cm, below=1.5cm of title.south, xshift=4.5cm] (h2)
  {ІОР оцінювання, контролю та рефлексії};
\node[dia/lbox=5.5cm, below=0.3cm of h2] (d2) {%
  -- засоби оцінювання\\
  -- хмаро орієнтовані тестові системи\\
  -- аналітика зворотного зв'язку%
};

% --- Lower-left ---
\node[dia/rbox=5.5cm, below=0.8cm of d1] (h3)
  {ІОР операційно-технологічної підготовки};
\node[dia/lbox=5.5cm, below=0.3cm of h3] (d3) {%
  -- практикуми\\
  -- тренажери\\
  -- навчально-методичні комплекси%
};

% --- Lower-right ---
\node[dia/rbox=5.5cm, below=0.8cm of d2] (h4)
  {ІОР комунікації та колаборативної взаємодії};
\node[dia/lbox=5.5cm, below=0.3cm of h4] (d4) {%
  -- комунікаційні засоби\\
  -- спільні імерсивні простори\\
  -- мережеві середовища взаємодії%
};

% --- Bottom-center ---
\coordinate (botcenter) at ($(d3.south)!0.5!(d4.south) + (0,-1.2cm)$);
\node[dia/rbox=7.5cm, anchor=north] at (botcenter) (h5)
  {ІОР педагогічного супроводу та методичної підтримки};
\node[dia/lbox=7.5cm, below=0.3cm of h5] (d5) {%
  -- програмно-методичні матеріали\\
  -- методичні рекомендації\\
  -- дидактичні демонстраційні матеріали\\
  -- посібники і підручники%
};

% Arrows: title -> upper headers
\draw[dia/arrow] (title.south) -- (h1.north);
\draw[dia/arrow] (title.south) -- (h2.north);

% Arrows: upper -> lower
\draw[dia/arrow] (d1.south) -- (h3.north);
\draw[dia/arrow] (d2.south) -- (h4.north);

% Arrows: lower -> bottom-center
\draw[dia/arrow] (d3.south) -- (h5.north);
\draw[dia/arrow] (d4.south) -- (h5.north);

\end{tikzpicture}%
%}

\caption{Функціональна класифікація імерсивних освітніх ресурсів у підготовці майбутніх учителів інформатики}
\label{fig:1.4}
\end{figure}

Класифікація відображає послідовність дидактичних функцій імерсивних освітніх ресурсів у процесі навчання майбутніх учителів інформатики, від проєктування та реалізації до рефлексивного оцінювання результатів.

Водночас ефективне проєктування ІОР потребує опори на моделі педагогічного дизайну, адаптовані до специфіки імерсивних технологій. Р.-Ю.~Лі та Дж.~Ч.-Ю.~Сан~\cite{Li2024ADDIE} на основі дослідження 23 ключових факторів обґрунтували динамічну модель ADDIE для проєктування VR-навчання, визначивши ``дизайн інтерактивного зворотного зв'язку'' як найвпливовіший фактор. П'ятифазна структура ADDIE (Аналіз $\to$ Дизайн $\to$ Розробка $\to$ Впровадження $\to$ Оцінювання) забезпечує системний підхід до створення ІОР. Альтернативний підхід запропоновано Б.~Черкавскі та М.~Берті~\cite{Czerkawski2021LXD}, які застосували модель послідовних наближень SAM~I до проєктування навчального досвіду з доповненою реальністю. Ітеративне швидке прототипування, притаманне SAM, виявляється більш придатним для розробки імерсивних ресурсів, де технічні обмеження вимагають постійного тестування.

У контексті цього дослідження принципово важливою є авторська 6-компонентна модель проєктування ІОР, обґрунтована С.~О.~Семеріковим, Т.~А.~Вакалюк, І.~С.~Мінтій та ін. за участі автора~\cite{Semerikov2022IERDesign}: Аналіз $\to$ Вимоги $\to$ Добір програмного забезпечення $\to$ Планування та розробка $\to$ Впровадження $\to$ Оцінювання. На відміну від універсальних моделей ADDIE та SAM, ця модель спеціально розроблена для ІОР і включає етап добору програмного інструментарію як окремий структурний компонент, що відповідає специфіці імерсивних технологій. Порівняння зазначених моделей подано у табл.~\ref{tab:design_models}.

\begin{table}
\caption{Порівняння моделей педагогічного дизайну для ІОР}\label{tab:design_models}\small
\begin{tabularx}{\linewidth}{|p{0.15\linewidth}|p{0.16\linewidth}|X|p{0.25\linewidth}|c|}
\hline
\hfil \textbf{Модель} & \hfil \textbf{Фази} & \textbf{Ключова особливість} & \hfil \textbf{Придатність для ІОР} & \textbf{Джерело} \\
\hline
ADDIE & 5, лінійних & Системність, структурованість & Планування, але обмежена ітеративність & \cite{Li2024ADDIE} \\
\hline
SAM~I & 3, ітеративних & Швидке прототипування & Технічна розробка AR/VR & \cite{Czerkawski2021LXD} \\
\hline
Авторська 6-компо\-нентна & 6, зі зворотним зв'язком & Інтегрує добір ПЗ & Спеціально для ІОР & \cite{Semerikov2022IERDesign} \\
\hline
LXD для AR & Ітеративна + поріг & Орієнтація на досвід учня & AR-специфічна & \cite{Czerkawski2021LXD} \\
\hline
\end{tabularx}
	
\end{table}


Авторська модель також передбачає трьохетапний підхід до навчання: теоретичний $\to$ практичний $\to$ апробаційний~\cite{Semerikov2022IERDesign}, що безпосередньо відповідає структурно-функціональній моделі ПНР~ІОР, представленій далі.

\emph{\textbf{Процес навчання розробки імерсивних освітніх ресурсів}} розглядається як цілісно організована педагогічна діяльність, спрямована на формування у майбутніх учителів інформатики системи знань, умінь і практичних навичок, необхідних для проєктування, створення та впровадження ефективних імерсивних освітніх ресурсів у навчальний процес. Зміст цього процесу ґрунтується на поєднанні \emph{теоретичної підготовки} з \emph{практично орієнтованою діяльністю} та охоплює опанування концептуальних засад імерсивних технологій, зокрема принципів функціонування віртуальної, доповненої та змішаної реальності, а також інших цифрових середовищ занурення. Важливим складником ПНР~ІОР є \emph{освоєння інструментів розробки} імерсивних ресурсів, що передбачає набуття досвіду роботи з програмними середовищами та платформами для створення VR-, AR- та MR-застосунків, зокрема Unity, Unreal Engine, ARKit, ARCore. Поряд із технічною підготовкою особливе значення надається \emph{педагогічному проєктуванню}, у межах якого майбутні вчителі інформатики навчаються розробляти дидактично обґрунтовані сценарії навчання, визначати освітні цілі та завдання, добирати адекватні методи й прийоми навчальної діяльності з урахуванням можливостей імерсивних середовищ.

Процес навчання також передбачає апробацію створених ІОР шляхом їх тестування та оцінювання, аналізу результатів використання й збирання зворотного зв'язку, що дозволяє коригувати зміст і функціональні характеристики цифрових продуктів. Завершальним компонентом ПНР~ІОР є розробка навчально-методичного забезпечення, яке охоплює створення інструктивних матеріалів, методичних рекомендацій та супровідної документації для ефективного використання імерсивних освітніх ресурсів у практиці закладів освіти. У сукупності зазначені компоненти забезпечують формування у майбутніх учителів інформатики професійних компетентностей, необхідних для самостійної розробки, педагогічно вмотивованого впровадження та оцінювання імерсивних освітніх ресурсів у сучасному освітньому середовищі.

Отже, на цьому етапі ми припустили можливість структурування моделі процесу навчання розробки імерсивних освітніх ресурсів (рис.~\ref{fig:1.5}).

\begin{figure}[!h]
\centering
%\resizebox{\textwidth}{!}{%
\begin{tikzpicture}[node distance=0.6cm]

% --- Section 1: МЕТА ---
\node[dia/header=14cm] (h1) {МЕТА};
\node[dia/lbox=14cm, below=0pt of h1] (b1) {%
формування готовності майбутніх учителів інформатики до педагогічно обґрунтованої розробки та впровадження ІОР%
};

% Arrow 1->2
\node[below=0.6cm of b1] (arr1) {};
\draw[dia/arrow] (b1.south) -- ++(0,-0.6cm);

% --- Section 2: Зміст ---
\node[dia/header=14cm, below=1.0cm of b1] (h2) {Зміст};
\node[dia/lbox=14cm, below=0pt of h2] (b2) {%
\checkmark{} основи імерсивних технологій (VR, AR, MR, 360°, симуляції);\\[2pt]
\checkmark{} інструменти та середовища розробки імерсивного контенту;\\[2pt]
\checkmark{} педагогічне проєктування ІОР;\\[2pt]
\checkmark{} методика використання ІОР в освітньому процесі;%
};

% Arrow 2->3
\draw[dia/arrow] (b2.south) -- ++(0,-0.6cm);

% --- Section 3: Етапи реалізації ---
\node[dia/header=14cm, below=1.0cm of b2] (h3) {Етапи реалізації};
\node[dia/lbox=14cm, below=0pt of h3] (b3) {%
\checkmark{} теоретико-орієнтаційний етап;\\[2pt]
\checkmark{} інструментально-практичний етап;\\[2pt]
\checkmark{} проєктувально-конструкторський етап;\\[2pt]
\checkmark{} апробаційно-рефлексивний етап;\\[2pt]
\checkmark{} методично-оформлювальний етап;%
};

% Arrow 3->4
\draw[dia/arrow] (b3.south) -- ++(0,-0.6cm);

% --- Section 4: Результат ---
\node[dia/header=14cm, below=1.0cm of b3] (h4) {Результат};
\node[dia/lbox=14cm, below=0pt of h4] (b4) {%
\checkmark{} сформовані професійні компетентності з розробки ІОР;\\[2pt]
\checkmark{} здатність до самостійного створення та впровадження ІОР;\\[2pt]
\checkmark{} готовність до інтеграції імерсивних технологій у навчальний процес.%
};

\end{tikzpicture}%


\caption{Структурно-функціональна модель процесу навчання розробки імерсивних освітніх ресурсів}
\label{fig:1.5}
\end{figure}

\emph{\textbf{Цифровий дизайн}} (далі - ЦД) в освіті та його диференційовані варіанти (digital instructional design, digital learning design, digital media design). Зазначимо, що у вітчизняному науковому дискурсі цифровий дизайн (Digital Design) тлумачиться переважно як ``творчий метод, процес і результат художньо-технічного проектування за допомогою інформаційних технологій промислових виробів та різних комплексів і систем, орієнтованих на найбільшу відповідність створюваних об'єктів потребам сучасної людини та навколишньому природному середовищу як з прагматичної так і з естетичної точки зору'' \cite[c.~11]{Ziaziun2000Philosophy}. Таке розуміння зосереджується передусім на візуально-технічному вимірі цифрових продуктів і лише опосередковано враховує їхній педагогічний потенціал.

Водночас у науково-педагогічному контексті функціонують близькі, але змістово відмінні варіанти: \emph{Digital instructional design} використовується для позначення проєктування навчального контенту з орієнтацією на результати навчання, структурування змісту та дидактичну доцільність; \emph{Digital learning design} пов'язується з конструюванням цифрового освітнього середовища як цілісної системи навчальних активностей, взаємодії та підтримки здобувачів освіти; \emph{Digital media design} акцентує увагу на візуально-комунікаційному аспекті цифрових ресурсів, мультимедійності та користувацькому досвіді.

Для цифрового дизайну ІОР усі три варіанти є релевантними та взаємодоповнюючими: інструктивний дизайн визначає дидактичну логіку імерсивного ресурсу; дизайн навчання забезпечує цілісність освітнього середовища; медійний дизайн відповідає за якість візуально-інтерактивного досвіду. М.~Мюлдерс, Й.~Бухнер та М.~Керрес~\cite{Mulders2020Framework} обґрунтовують, що проєктування імерсивних ресурсів виходить за межі традиційного мультимедійного дизайну, оскільки принципи когнітивної теорії мультимедійного навчання не можуть бути безпосередньо перенесені на іVR/AR-середовища з їх просторовою взаємодією та тілесним зануренням. У контексті моделі TPACK~\cite{BeldaMedina2022TPACK}, проаналізованої у підрозділі~\ref{subsec:1.1}, цифровий дизайн ІОР розміщується на перетині технологічних знань (\emph{TK}) та педагогічних знань (\emph{PK}), утворюючи технолого-педагогічні знання (\emph{TPK}), дефіцит яких у майбутніх учителів було констатовано в низці досліджень.

Специфіка цифрового дизайну WebAR-ресурсів полягає у поєднанні трьох вимірів: веб-розробка (HTML, CSS, JavaScript), тривимірне моделювання (створення та оптимізація GLTF/GLB-моделей) і AR-композиція (просторове розміщення, трекінг, інтерактивна логіка). Така унікальна комбінація не охоплюється жодною із традиційних моделей інструктивного дизайну, що додатково обґрунтовує потребу в авторській 6-компонентній моделі (табл.~\ref{tab:design_models}), яка інтегрує етап добору програмного інструментарію.

Відтак цифровий дизайн у контексті дослідження виконує функцію перехідної ланки між технологічними можливостями імерсивних технологій та їх педагогічно доцільним використанням у навчальному процесі, забезпечуючи трансформацію технічних рішень у дидактично вмотивовані освітні продукти.

Теоретичним підґрунтям тези про необхідність навчання майбутніх учителів саме \emph{розробці} ІОР, а не лише їх використанню, слугує конструкціоністська парадигма. Відповідно до положень конструкціонізму, навчання є найбільш ефективним, коли суб'єкт створює зовнішні артефакти, матеріалізуючи власне розуміння. У контексті імерсивних технологій цей принцип набуває особливої актуальності: Х.~Лі та Й.~Хван~\cite{Lee2022VRMaking} експериментально підтвердили, що майбутні вчителі, які створювали VR-навчальний контент для підручників K-12 (\emph{VR-Making}), продемонстрували статистично значуще покращення технологічної готовності, розвиток компетентностей 4C (критичне мислення, креативність, колаборація, комунікація) та підвищення сприйнятої педагогічної цінності імерсивних технологій. Отже, VR-Making операціоналізує конструкціоністський підхід для імерсивної технологічної освіти. Т.~Кі, Х.~Чжан та Р.~Б.~Кінг~\cite{Kee2023Experiential} доповнюють цю позицію з перспективи експерієнціального навчання, демонструючи, що AR/VR-середовища конструктивно вирівнюють стадії рефлексивного спостереження та активного експериментування, однак потребують явного педагогічного скафолдингу для стадії абстрактної концептуалізації.

На перетині цифрового дизайну, педагогічного проєктування та методики навчання інформатики формується \emph{\textbf{професійна готовність майбутнього вчителя інформатики до розробки ІОР}}\emph{\textbf{.}} Зміст професійної готовності охоплює поєднання теоретичних знань про засади імерсивних технологій, принципи педагогічного проєктування, основи цифрового дизайну та методику навчання інформатики з практичними вміннями проєктувати навчальні сценарії, використовувати інструменти розробки імерсивних ресурсів і створювати інтерактивні навчальні середовища, орієнтовані на активну участь здобувачів освіти. Водночас важливою складовою цієї готовності є особистісні якості майбутнього вчителя інформатики, зокрема креативність, ініціативність, відповідальність, здатність до самостійної діяльності та професійної співпраці, які забезпечують продуктивну роботу в умовах швидкої цифрової трансформації освіти.

Окремого значення набуває мотиваційний аспект професійної готовності, що виявляється у стійкій зацікавленості майбутнього вчителя інформатики у використанні імерсивних технологій в освітньому процесі, прагненні до постійного професійного розвитку та вдосконалення власної педагогічної практики. Саме мотиваційна складова визначає не лише готовність до опанування нових технологій, а й здатність критично осмислювати їх педагогічний потенціал та межі доцільності застосування.

\emph{Професійна готовність майбутнього вчителя інформатики до розробки імерсивних освітніх ресурсів} у межах цього дослідження розглядається як інтегративна характеристика особистості, що виявляється у здатності педагогічно обґрунтовано проєктувати, створювати та впроваджувати імерсивні освітні ресурси на основі поєднання теоретичних знань з імерсивних технологій, цифрового дизайну й методики навчання інформатики, практичних умінь роботи з відповідними інструментами та платформами, а також сформованих особистісно-мотиваційних якостей, необхідних для професійної діяльності в умовах цифрової трансформації освіти. Таке розуміння професійної готовності узгоджується з положеннями стандарту підготовки вчителя інформатики, відповідно до яких фахова діяльність педагога передбачає не лише володіння цифровими технологіями, а й здатність до їх педагогічно доцільного використання для досягнення результатів навчання, проєктування освітнього середовища та забезпечення активної пізнавальної діяльності здобувачів освіти. У термінах моделі TPACK~\cite{BeldaMedina2022TPACK} чотирикомпонентна структура професійної готовності (рис.~\ref{fig:1.6}) відповідає перетину технологічних (TK), педагогічних (PK) та предметних (CK) знань: мотиваційний компонент забезпечує готовність до інтеграції технологій у педагогічну практику; теоретичний -- формує предметно-технологічні знання (TCK); практичний -- розвиває технолого-педагогічні знання (TPK); а особистісний -- забезпечує рефлексивну здатність до синтезу всіх трьох вимірів (TPACK). Подальший розвиток AR, VR і MR зумовлює зростання їхнього значення у формуванні цифрової та проєктувальної компетентності студентів, а також посилює інтеграцію імерсивних рішень у професійну діяльність сучасного вчителя. У цьому контексті цифровий дизайн імерсивних середовищ постає як складник професійної діяльності вчителя інформатики, що поєднує технологічну, дидактичну та методичну компетентності.

\begin{figure}[!h]
\centering
%\resizebox{\textwidth}{!}{%
\begin{tikzpicture}

% Title
\node[dia/rbox=15cm] (title)
  {\textbf{Професійна готовність майбутнього вчителя інформатики \\до розробки ІОР}};

% Column headers — centers spaced 4cm apart: -6, -2, +2, +6
\node[dia/rbox=3.3cm, below=1.5cm of title.south, xshift=-6cm] (ch1)
  {\textbf{Мотиваційний акцент}};
\node[dia/rbox=3.3cm, below=1.5cm of title.south, xshift=-2cm] (ch2)
  {\textbf{Теоретичні знання}};
\node[dia/rbox=3.3cm, below=1.5cm of title.south, xshift=2cm] (ch3)
  {\textbf{Практичні вміння}};
\node[dia/rbox=3.3cm, below=1.5cm of title.south, xshift=6cm] (ch4)
  {\textbf{Особистісні якості}};

% Lines from title to column headers
\draw[dia/line] (title.south) -- (ch1.north);
\draw[dia/line] (title.south) -- (ch2.north);
\draw[dia/line] (title.south) -- (ch3.north);
\draw[dia/line] (title.south) -- (ch4.north);

% Content boxes
\node[dia/lrbox=3.3cm, below=0.5cm of ch1] (cb1) {%
професійний розвиток; готовність до опанування нових технологій; критичне мислення%
};

\node[dia/lrbox=3.3cm, below=0.5cm of ch2] (cb2) {%
основи імерсивних технологій; принципи педагогічного проєктування; цифровий дизайн; методика навчання інформатики%
};

\node[dia/lrbox=3.3cm, below=0.5cm of ch3] (cb3) {%
інструменти розробки; інтерактивне середовища%
};

\node[dia/lrbox=3.3cm, below=0.5cm of ch4] (cb4) {%
відповідаль\-ність; самостійна діяльність; професійна співпраця%
};

\end{tikzpicture}%
%}

\caption{Структурно-логічна схема професійної готовності майбутнього вчителя інформатики до розробки імерсивних освітніх ресурсів}
\label{fig:1.6}
\end{figure}

У сукупності зазначені на рис.~\ref{fig:1.6} компоненти професійної готовності утворюють підґрунтя для переходу до осмислення педагогічної ефективності імерсивних освітніх ресурсів, оскільки рівень такої ефективності безпосередньо залежить від здатності вчителя інтегрувати цифровий дизайн імерсивних середовищ у логіку навчального процесу та орієнтувати їх використання на досягнення визначених освітніх результатів.

\emph{Педагогічна ефективність імерсивних освітніх ресурсів} у межах цього дослідження розглядається як міра їх здатності забезпечувати досягнення визначених освітніх результатів за умови педагогічно доцільного проєктування та використання в навчальному процесі. На відміну від суто технологічної результативності, педагогічна ефективність ІОР зумовлюється не стільки рівнем їх технічної складності, скільки відповідністю дидактичним цілям, змісту навчання та особливостям пізнавальної діяльності здобувачів освіти.

У педагогічному вимірі ІОР набувають педагогічної цінності лише за умови їх інтеграції у логіку навчального процесу, що передбачає чітке визначення освітніх цілей, продумане сценарне проєктування навчальної діяльності, адекватний добір методів і форм роботи, а також рефлексивне оцінювання результатів навчання. Відтак педагогічна ефективність ІОР безпосередньо залежить від рівня професійної готовності вчителя інформатики, який виступає суб'єктом їх розробки, адаптації та впровадження.

Педагогічно ефективні ІОР сприяють активізації пізнавальної діяльності, формуванню стійкого навчального інтересу, розвитку практичних умінь і навичок, а також забезпечують контекстуальність навчання через моделювання реальних або наближених до реальних ситуацій. Водночас відсутність педагогічного проєктування або формальне використання ІТ зводить їх потенціал до рівня візуального доповнення навчального матеріалу, не забезпечуючи суттєвого впливу на якість освітніх результатів.

Отже, педагогічна ефективність імерсивних освітніх ресурсів постає не автономною характеристикою цифрового продукту, а похідною від цілісного процесу навчання їх розробки та впровадження. Саме у взаємозв'язку з професійною готовністю майбутнього вчителя інформатики вона розглядається як важливий критерій обґрунтування доцільності використання імерсивних технологій у сучасному освітньому середовищі.

Можемо констатувати, що педагогічна ефективність імерсивних освітніх ресурсів~-- це ступінь досягнення освітніх цілей за допомогою використання ІОР. Вона може бути оцінена за різними критеріями, які кореспондують з факторами моделі CAMIL~\cite{Makransky2021CAMIL}: \emph{засвоєння знань} (рівень розуміння та запам\textquotesingle ятовування навчального матеріалу~-- пов'язане з фактором когнітивного навантаження CAMIL); \emph{формування навичок} (розвиток практичних вмінь та навичок~-- пов'язане з факторами агентності та втілення); \emph{мотивація до навчання} (підвищення інтересу та активності~-- безпосередньо відповідає фактору мотивації CAMIL); \emph{залучення} (ефект присутності та залучення~-- кореспондує з фактором присутності); \emph{зручність використання} (простота та зручність ІОР~-- пов'язана з фактором саморегуляції, оскільки зручний інтерфейс знижує когнітивне навантаження на навігацію). Зазначені критерії узгоджуються також з результатами мета-аналізів, систематизованих у табл.~\ref{tab:meta_analysis} підрозділу~\ref{subsec:1.1}, де стійкий позитивний ефект імерсивних технологій (розміри ефекту g\,=\,0{,}82--d\,=\,0{,}98) зафіксовано саме за критеріями засвоєння знань і мотивації.

Водночас маємо зазначити: аналіз сучасних наукових джерел засвідчує (підрозділ~\ref{subsec:1.1}), що використання імерсивних технологій в освітньому процесі характеризується поєднанням виразного педагогічного потенціалу та низки суперечливих наслідків, які зумовлюють необхідність їх педагогічно виваженого впровадження. Дослідники наголошують, що імерсивні технології здатні істотно підвищувати якість навчання на різних рівнях освіти, насамперед за рахунок створення ефекту присутності та занурення, активізації пізнавальної діяльності й зростання навчальної мотивації \cite{Radianti2020VR}. Їх привабливість для здобувачів освіти пов'язується з можливістю інтерактивної взаємодії з навчальним контентом, розвитку творчого мислення та формування ключових компетентностей. Систематичний огляд Ван Цінь та Лі Яньпін засвідчує, що поєднання VR, AR та MR у педагогічній освіті створює нові можливості для адаптивного та інтерактивного навчання~\cite{Angelov2025Synergy}

Окремі дослідження підкреслюють, що імерсивне навчання трансформує традиційні підходи до засвоєння знань, перетворюючи навчальний процес на інтуїтивно зрозумілу, контекстуалізовану діяльність, у межах якої здобувачі освіти можуть експериментувати, здійснювати практичні дії та закріплювати знання у змодельованих ситуаціях \cite[c.~55--56]{Lytvynova2023Health}. Значущим аспектом позитивного впливу ІТ вважається їх потенціал у зниженні психологічної напруги, характерної для онлайн-навчання, а також у підтримці навчальної успішності в умовах цифровізації освітнього середовища \cite[c.~56]{Lytvynova2023Health}

Водночас системний аналіз літератури дозволяє виявити низку ризиків і проблем, що супроводжують упровадження імерсивних технологій. Серед них науковці виокремлюють неготовність здобувачів освіти, батьків і педагогів до педагогічно обґрунтованого використання комп'ютерно орієнтованих та імерсивних навчальних систем, що знижує ефективність їх застосування та може призводити до формалізації інновацій~\cite{Radianti2020VR}. Також наголошується на високій вартості обладнання й програмного забезпечення, складності технічного обслуговування та ризиках надмірної віртуалізації освітнього середовища, що потенційно обмежує живу комунікацію та соціальну взаємодію учасників освітнього процесу~\cite{DiNatale2020Immersive}

Зарубіжні вчені констатують існуванні двох протилежних позицій у науковому дискурсі: з одного боку, припущення про позитивний вплив IVR через підвищення залученості та емоцій, з іншого - гіпотеза про негативний вплив через когнітивне перевантаження та відволікання\cite{Parong2021Cognitive}. Отже окрему групу суперечливих наслідків становлять питання впливу імерсивних технологій на психофізичний стан здобувачів освіти, зокрема можливі негативні ефекти тривалого перебування у віртуальних середовищах, що розглядаються дослідниками як аспект, який потребує подальшого наукового вивчення та нормативного врегулювання~\cite{Parong2021Cognitive}

Узагальнено позиції вчених що впливу ІТ на освітній процес подано у табл.~\ref{tab:iot_effects}.

{\small
\begin{longtable}{|p{0.18\linewidth}|p{0.39\linewidth}|p{0.38\linewidth}|}
\caption{Позитивні та суперечливі наслідки використання ІТ в освіті}\label{tab:iot_effects}\\
\hline
\hfil \textbf{Аспект} & \hfil \textbf{Позитивні наслідки} & \hfil \textbf{Суперечливі наслідки / ризики} \\
\hline
\endfirsthead
\multicolumn{3}{l}{\textit{Продовження таблиці~\ref{tab:iot_effects}}}\\
\hline
\hfil \textbf{Аспект} & \hfil \textbf{Позитивні наслідки} & \hfil \textbf{Суперечливі наслідки / ризики} \\
\hline
\endhead
%\hline
\endfoot
Дидактичний & Поглиблене засвоєння матеріалу; контекстуальність навчання; відпрацювання практичних навичок у безпечному середовищі~\cite{Radianti2020VR, Dalgarno2010What} & Формалізація використання без педагогічного проєктування; перевантаження візуальними ефектами \\
\hline
Мотиваційний & Зростання навчальної мотивації; зацікавленість; розвиток творчого мислення~\cite{Angelov2025Synergy, KMotivation2023} & Залежність мотивації від технологічної новизни \\
\hline
Особистісний & Розвиток соціальних навичок, співпраці, комунікації у віртуальних середовищах~\cite{DiNatale2020Immersive} & Зниження рівня безпосередньої соціальної взаємодії \\
\hline
Організаційний & Розширення доступу до ресурсів; подолання географічних і матеріальних обмежень~\cite{Radianti2020VR} & Висока вартість обладнання; складність технічного супроводу \\
\hline
Професійно-педагогічний & Можливості модернізації підготовки майбутніх учителів; розвиток цифрових і педагогічних компетентностей~\cite{Angelov2025Synergy} & Неготовність педагогів; потреба у спеціальній підготовці \\
\hline
Здоров'язбе\-режувальний & Потенціал зниження психологічної напруги в онлайн-навчанні~\cite{Krokos2019Virtual} & Неоднозначний вплив на психофізичний стан здобувачів освіти~\cite{Parong2021Cognitive} \\
\hline
Когнітивний & AR знижує стороннє когнітивне навантаження, оптимізуючи навчання початківців~\cite{Poupard2024Systematic, Buchner2021CogLoad} & VR може збільшувати когнітивне навантаження; ефект залежить від дизайну~\cite{Parong2021Cognitive, Mulders2020Framework} \\
\hline
\end{longtable}
}

Подані у таблиці відомості засвідчують, що педагогічний потенціал імерсивних технологій реалізується у різних вимірах освітнього процесу~-- дидактичному, мотиваційному, особистісному, організаційному та професійно-педагогічному. Водночас кожен із позитивних ефектів супроводжується певними ризиками, що зумовлює необхідність системного та педагогічно виваженого підходу до впровадження імерсивних технологій в освіті. Така подвійність результатів підкреслює, що ефективність використання імерсивних технологій визначається не стільки рівнем їх технологічної складності, скільки умовами педагогічного проєктування, підготовленістю вчителя та методичною доцільністю інтеграції у навчальний процес.

Узагальнення результатів теоретичного аналізу дозволяє представити імерсивні технології в освіті як цілісне педагогічне явище, що поєднує технологічні, дидактичні та особистісно-професійні компоненти (рис.~\ref{fig:1.7}). У такому взаємозв'язку імерсивні технології постають не лише як сукупність цифрових інструментів, а як основа для проєктування імерсивних освітніх ресурсів, організації процесу навчання їх розробки, формування професійної готовності майбутнього вчителя інформатики та досягнення педагогічної ефективності навчання.

\begin{figure}[!h]
\centering
%\resizebox{\textwidth}{!}{%
\begin{tikzpicture}
\small

% Title
\node[dia/header=15cm] (title) {ІМЕРСИВНІ ТЕХНОЛОГІЇ В ОСВІТІ};

% Column headers — all at same y, minimum height ensures uniform height
\node[dia/rbox=3.0cm, minimum height=1.2cm, below=1.5cm of title.south, xshift=-5.7cm] (col1h) {\textbf{Імерсивні технології}};
\node[dia/rbox=3.0cm, minimum height=1.2cm, below=1.5cm of title.south, xshift=-1.9cm] (col2h) {\textbf{Імерсивні освітні ресурси}};
\node[dia/rbox=3.0cm, minimum height=1.2cm, below=1.5cm of title.south, xshift=1.9cm] (col3h) {\textbf{Цифровий дизайн в освіті}};
\node[dia/rbox=3.0cm, minimum height=1.2cm, below=1.5cm of title.south, xshift=5.7cm] (col4h) {\textbf{Педагогічна\\ ефективність ІОР}};

% Row a — all anchored to same y (relative to col1h bottom)
\node[dia/rbox=3.0cm, below=0.6cm of col1h] (col1a) {Віртуальна реальність};
\node[dia/rbox=3.0cm, anchor=center] at (col2h |- col1a) (col2a) {інтерактивність};
\node[dia/rbox=3.0cm, anchor=center] at (col3h |- col1a) (col3a) {інтерактивність};
\node[dia/rbox=3.0cm, anchor=center] at (col4h |- col1a) (col4a) {Засвоєння знань};

% Row b
\node[dia/rbox=3.0cm, below=0.6cm of col1a] (col1b) {Доповнена реальність};
\node[dia/rbox=3.0cm, anchor=center] at (col2h |- col1b) (col2b) {занурення};
\node[dia/rbox=3.0cm, anchor=center] at (col3h |- col1b) (col3b) {інтерактивність};
\node[dia/rbox=3.0cm, anchor=center] at (col4h |- col1b) (col4b) {Формування навичок};

% Row c
\node[dia/rbox=3.0cm, below=0.6cm of col1b] (col1c) {Змішана реальність};
\node[dia/rbox=3.0cm, anchor=center] at (col2h |- col1c) (col2c) {контекстуальність};
\node[dia/rbox=3.0cm, anchor=center] at (col3h |- col1c) (col3c) {інтерактивність};
\node[dia/rbox=3.0cm, anchor=center] at (col4h |- col1c) (col4c) {Мотивація до навчання};

% Vertical lines from title to column headers
\draw[dia/line] (title.south) -- ++(0,-0.5cm) -| (col1h.north);
\draw[dia/line] (title.south) -- ++(0,-0.5cm) -| (col2h.north);
\draw[dia/line] (title.south) -- ++(0,-0.5cm) -| (col3h.north);
\draw[dia/line] (title.south) -- ++(0,-0.5cm) -| (col4h.north);

% Horizontal arrows between sub-items
\draw[dia/arrow] (col1a.east) -- (col2a.west);
\draw[dia/arrow] (col2a.east) -- (col3a.west);
\draw[dia/arrow] (col3a.east) -- (col4a.west);

\draw[dia/arrow] (col1b.east) -- (col2b.west);
\draw[dia/arrow] (col2b.east) -- (col3b.west);
\draw[dia/arrow] (col3b.east) -- (col4b.west);

\draw[dia/arrow] (col1c.east) -- (col2c.west);
\draw[dia/arrow] (col2c.east) -- (col3c.west);
\draw[dia/arrow] (col3c.east) -- (col4c.west);

% Bottom section
\coordinate (botsep) at ($(col1c.south)!0.5!(col4c.south) + (0,-1.5cm)$);
\node[dia/header=15cm, anchor=north] at (botsep) (bottitle) {Процес навчання розробки ІОР};

% Three boxes in a chain
\node[dia/rbox=4cm, below=1.0cm of bottitle.south, xshift=-5.0cm] (step1)
  {Вивчення основ ІТ};
\node[dia/rbox=4cm, below=1.0cm of bottitle.south] (step2)
  {Освоєння інструментів розробки};
\node[dia/rbox=4cm, below=1.0cm of bottitle.south, xshift=5.0cm] (step3)
  {Педагогічне проєктування};

% Arrows between steps
\draw[dia/arrow] (step1.east) -- (step2.west);
\draw[dia/arrow] (step2.east) -- (step3.west);

% Lines from bottom title to steps
\draw[dia/line] (bottitle.south) -- ++(0,-0.3cm) -| (step1.north);
\draw[dia/line] (bottitle.south) -- ++(0,-0.3cm) -| (step2.north);
\draw[dia/line] (bottitle.south) -- ++(0,-0.3cm) -| (step3.north);

\end{tikzpicture}%
%}

\caption{Структурно-логічна схема імерсивних технологій, представлених в освіті}
\label{fig:1.7}
\end{figure}

Представлена структурно-логічна схема відображає взаємозв'язок технологічних, дидактичних і особистісно-професійних компонентів, що визначають потенціал імерсивних технологій у сучасному освітньому середовищі. Разом із тим теоретичне осмислення імерсивних технологій потребує доповнення аналізом практичних способів їх реалізації в освітньому процесі. Саме методичні підходи до використання імерсивних технологій визначають, яким чином технологічні можливості трансформуються у педагогічні рішення, що забезпечують досягнення результатів навчання, урахування індивідуальних освітніх потреб та підвищення якості освітнього процесу. З огляду на це доцільним є звернення до вітчизняного досвіду, у межах якого напрацьовано різноманітні методичні підходи до інтеграції імерсивних технологій у навчання. Їх вивчення дозволяє виявити усталені практики, інноваційні рішення та тенденції, що мають значення для подальшого обґрунтування методичних засад використання імерсивних технологій у вітчизняній освітній практиці, чому й присвячено наступний підрозділ.

